\section{Experimental Evaluation}
\label{sec:evaluation}

In this section, we evaluate the implemented static GCN baseline for build failure prediction. The goal is not to claim state-of-the-art performance, but to quantify what can and cannot be learned from whole-project static call graphs plus CI metadata.

\subsection{Research Questions}

Our evaluation is guided by three research questions. RQ1, predictive performance, asks whether a lightweight static GCN can predict failing builds better than simple baselines. RQ2, dataset sensitivity, asks how stable the results are across the available metadata datasets. RQ3, graph failure mode, asks whether the graph representation encodes useful build-specific dependency information or mostly behaves like a project identifier.

\subsection{Evaluation Protocol}

We evaluate three logical datasets: \texttt{dataset\_init.csv}, \texttt{dataset.csv}, and \texttt{dataset\_filter.csv}. The first file uses an older metadata schema, so we normalize it into the current training schema by mapping \texttt{ts\_duration}, \texttt{passcount}, and \texttt{failcount} into \texttt{test\_suite\_duration\_s}, \texttt{num\_pass\_class}, and \texttt{num\_fail\_class}; because it has no stage field, we assign \texttt{stage\_id=unknown}. For all datasets, the binary target is
\[
y_i=\mathbb{1}[\texttt{num\_fail\_class}_i>0].
\]

To approximate deployment in CI, build records are sorted by timestamp. The first 70\% are used for training, the next 15\% for validation, and the final 15\% for testing. The validation set is used to select the decision threshold $\tau$ that maximizes F1 over thresholds in $[0.2,0.8]$.

\subsection{Metrics and Baseline}

Let TP, TN, FP, and FN denote true positives, true negatives, false positives, and false negatives. We report
\[
\text{Accuracy}=\frac{\mathrm{TP}+\mathrm{TN}}{\mathrm{TP}+\mathrm{TN}+\mathrm{FP}+\mathrm{FN}},
\]
\[
\text{Precision}=\frac{\mathrm{TP}}{\mathrm{TP}+\mathrm{FP}},
\qquad
\text{Recall}=\frac{\mathrm{TP}}{\mathrm{TP}+\mathrm{FN}},
\]
\[
\text{F1}=\frac{2\cdot\text{Precision}\cdot\text{Recall}}{\text{Precision}+\text{Recall}}.
\]
These metrics are more informative than accuracy because the failure label is highly imbalanced and the positive rate changes over time.

We compare against an all-positive baseline that predicts every build as failing. If the positive rate is $\pi=N_+/N$, then this baseline has
\[
\text{Accuracy}=\pi,\qquad
\text{Precision}=\pi,\qquad
\text{Recall}=1,
\]
and therefore
\[
\text{F1}_{\mathrm{all+}}=\frac{2\pi}{1+\pi}.
\]
This baseline is intentionally simple, but it is mathematically important: when the positive rate is high, always predicting failure can achieve a strong F1 score.

\subsection{Overall Results}

\begin{table*}[ht]
    \centering
    \caption{Cross-dataset performance of the static GCN build failure predictor.}
    \label{tab:cross_dataset_results}
    \begin{tabular}{llrrrrrrrr}
        \toprule
        Dataset & Split & Rows & Pos. Rate & Best Epoch & $\tau$ & Acc. & Prec. & Rec. & F1 \\
        \midrule
        \texttt{dataset\_init} & Train & 2,334 & 0.891 & 40 & 0.20 & 0.919 & 0.973 & 0.934 & 0.953 \\
        \texttt{dataset\_init} & Validation & 500 & 0.434 & 40 & 0.20 & 0.732 & 0.799 & 0.512 & 0.624 \\
        \texttt{dataset\_init} & Test & 501 & 0.425 & 40 & 0.20 & 0.585 & 0.508 & 0.728 & 0.598 \\
        \midrule
        \texttt{dataset} & Train & 5,157 & 0.902 & 34 & 0.20 & 0.926 & 0.963 & 0.954 & 0.959 \\
        \texttt{dataset} & Validation & 1,105 & 0.702 & 34 & 0.20 & 0.837 & 0.906 & 0.857 & 0.881 \\
        \texttt{dataset} & Test & 1,106 & 0.389 & 34 & 0.20 & 0.741 & 0.961 & 0.347 & 0.509 \\
        \midrule
        \texttt{dataset\_filter} & Train & 5,157 & 0.902 & 34 & 0.20 & 0.926 & 0.963 & 0.954 & 0.959 \\
        \texttt{dataset\_filter} & Validation & 1,105 & 0.702 & 34 & 0.20 & 0.837 & 0.906 & 0.857 & 0.881 \\
        \texttt{dataset\_filter} & Test & 1,106 & 0.389 & 34 & 0.20 & 0.741 & 0.961 & 0.347 & 0.509 \\
        \bottomrule
    \end{tabular}
\end{table*}

Table~\ref{tab:cross_dataset_results} answers RQ1 and RQ2. The model fits the training data and performs well on validation for the current datasets, but it does not generalize reliably to the final temporal test period. On \texttt{dataset} and \texttt{dataset\_filter}, test precision is high at 0.961, but recall falls to 0.347. In other words, when the model predicts failure it is usually correct, but it misses most actual failures.

\begin{table}[ht]
    \centering
    \caption{Test F1 compared with the all-positive baseline.}
    \label{tab:baseline_results}
    \begin{tabular}{lrr}
        \toprule
        Dataset & All-positive F1 & GCN F1 \\
        \midrule
        \texttt{dataset\_init} & 0.597 & 0.598 \\
        \texttt{dataset} & 0.560 & 0.509 \\
        \texttt{dataset\_filter} & 0.560 & 0.509 \\
        \bottomrule
    \end{tabular}
\end{table}

Table~\ref{tab:baseline_results} shows that the GCN is effectively tied with the all-positive baseline on \texttt{dataset\_init} and worse than it on the two current datasets. This is a strong negative result: a graph-based method should at least outperform a classifier with no graph information.

\subsection{Temporal and Project Distribution Shift}

The temporal split creates a substantial distribution shift. In \texttt{dataset} and \texttt{dataset\_filter}, the training period is dominated by \texttt{kafka}: 4,589 of 5,157 training rows, with a 0.964 positive rate. The test period contains many more \texttt{james}, \texttt{hadoop}, and \texttt{hbase} rows, with much lower or more mixed positive rates. For example, \texttt{james} contributes 422 test rows but has only a 0.052 positive rate, while \texttt{kafka} contributes 137 test rows with a 0.993 positive rate.

This matters because the graph embedding $g_p$ is fixed per project. If training associates $g_{\texttt{kafka}}$ with frequent failures and associates other project embeddings with fewer observed failures, the classifier can appear strong during validation but fail when the project mix changes.

\subsection{Per-Project Test Behavior}

\begin{table*}[ht]
    \centering
    \caption{Per-project test metrics for \texttt{dataset} and \texttt{dataset\_filter}.}
    \label{tab:per_project_results}
    \begin{tabular}{lrrrrrrrrr}
        \toprule
        Project & Rows & Pos. Rate & TP & TN & FP & FN & Prec. & Rec. & F1 \\
        \midrule
        \texttt{hadoop} & 287 & 0.467 & 6 & 152 & 1 & 128 & 0.857 & 0.045 & 0.085 \\
        \texttt{hbase} & 257 & 0.537 & 6 & 118 & 1 & 132 & 0.857 & 0.043 & 0.083 \\
        \texttt{james} & 422 & 0.052 & 1 & 397 & 3 & 21 & 0.250 & 0.045 & 0.077 \\
        \texttt{kafka} & 137 & 0.993 & 136 & 0 & 1 & 0 & 0.993 & 1.000 & 0.996 \\
        \texttt{karaf} & 3 & 0.000 & 0 & 3 & 0 & 0 & 0.000 & 0.000 & 0.000 \\
        \bottomrule
    \end{tabular}
\end{table*}

Table~\ref{tab:per_project_results} answers RQ3. The model almost perfectly captures \texttt{kafka} failures but misses nearly all failures in \texttt{hadoop}, \texttt{hbase}, and \texttt{james}. This behavior is consistent with a project-prior model, not a build-specific dependency model.

\subsection{Graph Structure Diagnostics}

\begin{table*}[ht]
    \centering
    \caption{Static call graph statistics. The largest component is reported as a percentage of nodes.}
    \label{tab:graph_stats}
    \begin{tabular}{lrrrrrr}
        \toprule
        Project & Nodes & Edges & Avg. Degree & Weak Comp. & Largest Comp. & Max Degree \\
        \midrule
        \texttt{hadoop} & 87,016 & 253,694 & 5.83 & 974 & 96.9\% & 5,828 \\
        \texttt{hbase} & 42,609 & 131,416 & 6.17 & 470 & 96.2\% & 3,207 \\
        \texttt{jackrabbit-oak} & 34,245 & 94,337 & 5.51 & 426 & 96.1\% & 2,122 \\
        \texttt{kafka} & 29,757 & 87,855 & 5.90 & 577 & 95.0\% & 2,889 \\
        \texttt{james} & 23,723 & 69,061 & 5.82 & 337 & 95.9\% & 2,062 \\
        \texttt{karaf} & 7,870 & 21,794 & 5.54 & 166 & 94.3\% & 597 \\
        \bottomrule
    \end{tabular}
\end{table*}

The graph statistics in Table~\ref{tab:graph_stats} explain why the graph signal is weak. These are large sparse graphs with one giant weak component and many small components. Since the model mean-pools over all nodes, localized information is diluted:
\[
g_p=\frac{1}{|V_p|}\sum_{v\in V_p}H_{p,v}^{(2)}.
\]
If a failure-relevant region contains $k$ nodes in a graph with $|V_p|$ nodes, its direct contribution to the mean is only $k/|V_p|$. For a project such as \texttt{hadoop}, even a 100-node affected region contributes approximately
\[
\frac{100}{87016}\approx 0.00115
\]
of the unweighted mean. Without changed-node indicators or coverage-guided pooling, the model has little chance to preserve this local signal.

\subsection{Summary of Findings}

The static GCN fails for three mathematically concrete reasons. First, $g_p$ is invariant across builds from the same project, so the graph cannot represent build-specific changes. Second, symmetrization replaces $A_p$ with $A_p+A_p^\top+I$, discarding the direction of calls. Third, mean pooling over large graphs compresses tens of thousands of node embeddings into one vector, which washes out localized fault information. These limitations explain why the model performs like a project-prior classifier rather than a dependency-aware failure predictor.
